\chapter{Accuracy-driven certification}\label{ch:accuracy}
\chapterlead{A requested tolerance is an input to the problem, not an informal expectation
about its answer. The controller accepts a result only when an outward
probability interval proves the multiplicative accuracy requirement. Failure
to prove that inequality remains a failure of certification, even when a
point estimate appears accurate.}
\label{sec:final-5}

The public numerical contract should specify the desired accuracy, not the
internal transform resolution.

For an exact positive spectrum and threshold \(q\), the solver accepts a
requested relative multiplicative accuracy
\[
\varepsilon_{\rm req}>0
\]
and returns a certified interval
\[
[p_-,p_+]
\]
only when
\[
\boxed{
\frac{p_+}{p_-}-1
\le
\varepsilon_{\rm req}.
}
\]

In logarithmic coordinates, if
\[
\log_{10}p\in[d_-,d_+],
\]
the acceptance rule is
\[
\boxed{
10^{d_+-d_-}-1
\le
\varepsilon_{\rm req}.
}
\]

No value of \(N\), support width, or arithmetic precision appears in the
external contract.

\begin{theorem}[Accuracy contract for the positive exact-spectrum backend]\label{claim:final-5-1}
Assume an exact supplied positive decimal spectrum of rank greater than two,
a threshold \(q\ge0\) with valid interval-design parameters, and requested multiplicative tolerance
\(\varepsilon_{\rm req}>0\).  Suppose the directed-rounding interval backend
returns an enclosure
\[
\log_{10}p\in[d_-,d_+]
\]
that includes the proved support/alias and omitted-mode remainders.  If the
accuracy-driven controller returns \texttt{certified=True}, then
\[
\boxed{
P(Q\ge q)\in[10^{d_-},10^{d_+}]
}
\]
and
\[
\boxed{
\frac{10^{d_+}}{10^{d_-}}-1
\le
\varepsilon_{\rm req}.
}
\]
If the resource caps are reached before this inequality is proved, the
controller returns \texttt{certified=False} rather than promoting the interval
midpoint to a certified answer.
\end{theorem}

\begin{proof}
The directed-rounding backend encloses the finite transform and exponential
prefactor and then enlarges that enclosure by the proved support/alias and
infinite-mode error bounds.  Therefore the true tail probability lies inside
the returned interval.  The controller encloses
\(\operatorname{expm1}(\log(10)(d_+-d_-))\) by directed rounding and
sets \texttt{certified=True} only when the upper endpoint of that ratio
enclosure is no larger than the lower endpoint of the exact requested
tolerance enclosure.  High precision alone is not an enclosure proof.  The fail-closed branch follows directly from
the controller definition.
\end{proof}

\section{Adaptive refinement}

The controller adjusts three internal resources:

\begin{enumerate}
\item the COS mode count \(N\).
\item the tilted support scale \(L\), which changes \(b\) and its associated
Chernoff parameter.
\item the interval arithmetic precision.
\end{enumerate}

At each stage the full directed-rounding interval is recomputed.  If the
rounding component is material relative to the requested budget, arithmetic
precision is increased.  Otherwise the dominant analytic contribution
determines the next refinement: mode error increases \(N\), while
support/alias error increases \(L\).

A certified return occurs only after the multiplicative inequality above
has been proved. Reaching a resource limit instead terminates the calculation
with the failure status described next.

\section{Fail-closed behavior}

If the configured resource caps are exhausted before the requested accuracy
is established, the solver does not silently return its midpoint as though it
were certified.

Instead it returns
\[
\boxed{\texttt{certified = false}}
\]
together with the unresolved interval and a failure reason.

This behavior is part of the numerical definition of the solver.

\section{Accuracy ladder through \texorpdfstring{\(10^{-5000}\)}{1e-5000}}

The exact control
\[
Q=0.02\chi^2_{80}
\]
was evaluated at tail depths
\[
10^{-100},
\qquad
10^{-1000},
\qquad
10^{-5000}.
\]

For each depth the controller was asked to prove relative multiplicative
accuracies
\[
10^{-8},
\quad
10^{-10},
\quad
10^{-12},
\quad
10^{-14}.
\]
All twelve contracts were satisfied.

At the \(10^{-5000}\) point, the controller produced:

\begin{center}
\begin{tabular}{@{}rrrr@{}}
\toprule
requested accuracy & achieved ratio minus one & \(N\) & \(L\)\\
\midrule
\(10^{-8}\)  & \(3.19\times10^{-10}\) & 64  & 10\\
\(10^{-10}\) & \(7.70\times10^{-11}\) & 128 & 10\\
\(10^{-12}\) & \(3.48\times10^{-15}\) & 128 & 12\\
\(10^{-14}\) & \(3.48\times10^{-15}\) & 128 & 12\\
\bottomrule
\end{tabular}
\end{center}

The gap between requested and achieved accuracy is expected: refinement is
discrete, so one additional support or mode refinement can improve the
certificate by several orders of magnitude.

A further stress test requested
\[
10^{-16}
\quad\text{and}\quad
10^{-18}
\]
relative accuracy at the same \(10^{-5000}\) tail.  The solver obtained
approximately
\[
3.05\times10^{-18}
\]
and
\[
9.20\times10^{-20}
\]
respectively.

\begin{figure}[htbp]
\centering
\includegraphics[width=\linewidth]{figures_production/figure_26_accuracy_driven_refinement.pdf}
\caption{Requested and achieved multiplicative widths, expressed as decimal accuracy orders. Larger values mean tighter widths, so points above the equality line meet the requested tolerance. The two deepest stress results use the rounded numerical summaries in the text. Acceptance itself uses the full directed endpoints.}
\label{fig:accuracy-driven}
\end{figure}

\section{Acceptance arithmetic and serialization are directed}

The acceptance condition must be evaluated as carefully as the probability
interval.

Converting a tiny log-width \(w\) using
\[
10^w-1
\]
under an insufficient point-precision context can lose the width. At approximately \(10^{-15}\)
relative scales, the ratio can round to zero.

Even when the probability interval remains valid, that conversion cannot
support a certification decision. The implementation encloses
\[
\boxed{
\operatorname{expm1}(\log(10)\,w)
}
\]
with directed interval arithmetic before deciding whether the requested
accuracy has been met.  Acceptance compares the upper endpoint of that ratio
with the lower endpoint of the requested exact tolerance, so a rounded point
estimate cannot promote an unresolved boundary case to success.

Authoritative output strings must preserve the same outer enclosure.  The
implementation converts each exact binary interval endpoint to a decimal
string with \texttt{Decimal} and explicit \texttt{ROUND\_FLOOR} for lower
endpoints or \texttt{ROUND\_CEILING} for upper endpoints.  Ordinary interval
display formatting is not used as an authoritative serialization method:
nearest-rounded decimal text can move an endpoint inward even when the
internal interval arithmetic was directed.  Matrix-to-tail propagation keeps
these endpoint strings and applies the same directed ratio test.

The correctness requirement applies not only to the tail integral but also
to the logic that labels its result certified.

\section{Input semantics at \texorpdfstring{\(p\approx10^{-16}\)}{p about 1e-16}}

At this accuracy scale, the representation of the threshold \(q\) itself is
measurable.  In the \(0.02\chi^2_{80}\) control targeted at \(10^{-16}\),
rounding the high-precision threshold to binary float64 shifts the exact tail
by approximately
\[
3.83\times10^{-15}
\]
decades.  The hardened float64 solver evaluated at that same rounded input
then differs from the exact tail of the rounded input by approximately
\[
5.17\times10^{-15}
\]
decades.

Accordingly, claims below the float64 input-quantization scale use exact-decimal
or interval inputs.  The rigorous accuracy contract therefore certifies the
mathematical inputs supplied to it, rather than silently identifying a rounded
binary threshold with an ideal decimal threshold.



